% ============================================
% CONT-10
% Gerado de molambudos.md — build_miolo.py
% ============================================
\section*{\textbf{CONT-10} --- A Ferida}
\markboth{CONT-10 --- A Ferida}{}

\noindent \textit{A mancha na pele}

\noindent Você sentiu?

\noindent No meio da leitura, uma coceira no braço direito. Você passou a mão. Depois coçou. Depois olhou.

\noindent Há uma mancha na sua pele.

\noindent Pequena. Quase imperceptível. Mas está lá. Um ponto escuro, como tinta que vazou por baixo da pele.

\noindent Você passa o dedo. A mancha não sai. Passa com saliva. Não sai.

\noindent Sinta: seu coração bate mais forte bem naquele ponto da pele, como se a mancha tivesse pulso próprio.

\noindent A mancha é embaixo da pele.

\noindent Ela está \textit{crescendo}.

\noindent Não agora. Nas próximas horas. Amanhã de manhã, ela vai estar maior. Depois de amanhã, vai ter forma.

\noindent Forma de letra.

\noindent Uma letra torta. Infantil.

\noindent A primeira letra do nome de Joaquim.

\noindent \textit{\textquotedblleft{}J\textquotedblright{}.}

\noindent Amanhã: \textit{\textquotedblleft{}J\textquotedblright{}}.
Depois: \textit{\textquotedblleft{}JO\textquotedblright{}}.
Dia seguinte: \textit{\textquotedblleft{}JOA\textquotedblright{}}.

\noindent Até formar: \textit{\textquotedblleft{}JOAQUIM\textquotedblright{}}.

\noindent Na sua pele.

\noindent Como uma tatuagem que ninguém fez.
Como uma marca que veio de dentro.

\noindent Você pensa em mostrar a alguém. Um médico. Alguém da família. Mas o que diria? \textquotedblleft{}Tem uma letra crescendo debaixo da minha pele\textquotedblright{}? Você já sabe como isso soaria. Então escolhe não mostrar a ninguém. Escolhe ficar com o medo sozinho, do jeito que Joaquim ficou.

\noindent Olhe de novo.

\noindent A criatura está se revelando.
Não no livro.
Em você.

\vspace{0.5em}

\noindent\textit{\textbf{↪ Rotas:} \rota{CONT-11} \rotasep \rota{CONT-07}}